% ============================================================
% Formulario, argomento 1: eventi e insiemi (dispensa, capitolo 1)
% ============================================================
\argomento{Eventi e insiemi}{Dispensa: capitolo 1}

\voce{operazioni}{Operazioni fra eventi}{il vocabolario degli eventi}
  {A \cup B \ \text{(almeno uno)} \qquad A \cap B \ \text{(entrambi)} \qquad \compl{A} \ \text{(non } A\text{)}
   \qquad A \setminus B \ \text{(}A \text{ ma non } B\text{)}}
  {\variante{incompatibili}{A \cap B = \emptyset}
   \variante{$A$ implica $B$}{A \subseteq B}
   \variante{evento certo, evento impossibile}{\Omega, \quad \emptyset}}

\voce{differenza}{Differenza come intersezione}{}
  {A \setminus B = A \cap \compl{B}}
  {}

\voce{complementare}{Proprietà del complementare}{}
  {\compl{\compl{A}} = A \qquad A \cup \compl{A} = \Omega \qquad A \cap \compl{A} = \emptyset}
  {\variante{evento certo e impossibile}{\compl{\Omega} = \emptyset, \quad \compl{\emptyset} = \Omega}}

\voce{demorgan}{Leggi di De Morgan}{}
  {\compl{A \cup B} = \compl{A} \cap \compl{B} \qquad \compl{A \cap B} = \compl{A} \cup \compl{B}}
  {\variante{con $M$ eventi \oltre}{\compl{\bigcup_{i=1}^{M} A_i} = \bigcap_{i=1}^{M} \compl{A_i}, \quad
                                   \compl{\bigcap_{i=1}^{M} A_i} = \bigcup_{i=1}^{M} \compl{A_i}}}

\voce{associativa}{Proprietà associativa}{}
  {(A \cup B) \cup C = A \cup (B \cup C) \qquad (A \cap B) \cap C = A \cap (B \cap C)}
  {}

\voce{distributiva}{Proprietà distributive}{}
  {A \cap (B \cup C) = (A \cap B) \cup (A \cap C) \qquad A \cup (B \cap C) = (A \cup B) \cap (A \cup C)}
  {\variante{con $M$ eventi}{A_1 \cap \bigcup_{i=2}^{M} A_i = \bigcup_{i=2}^{M} (A_1 \cap A_i), \quad
                             A_1 \cup \bigcap_{i=2}^{M} A_i = \bigcap_{i=2}^{M} (A_1 \cup A_i)}}

\voce{cartesiano}{Prodotto cartesiano}{coppie ordinate}
  {A \times B = \{(a, b) : a \in A,\ b \in B\}}
  {\variante{un insieme per sé stesso}{A \times A = A^2, \quad A \times A \times A = A^3}
   \variante{numero di elementi (\vv{fondamentale})}{\card{A \times B} = \card{A} \cdot \card{B}}}
